% GENERATED FILE. Edit canonical lessons, then run npm run generate:books.
% canonical-source-hash: fnv1a64:605bebda50527c9b
% canonical-lessons: LA-C419-persolvo, LA-C419-usura, LA-C419-pignus, LA-C419-syngrapha, LA-C419-fenerator

\chapter{Paying It Back}
\label{ch:la-paying-it-back}
\hlchaptermodality{419}

\begin{chapteropening}
\textbf{By the end of this chapter:} \emph{I can say that a debt was paid in full, and talk about interest, a pledge, a written bond and a moneylender.}

Complete the last lesson of chapter 419: I can say that a debt was paid in full, and talk about interest, a pledge, a written bond and a moneylender.
\end{chapteropening}

\section[persolvō, persolvere]{persolvō, persolvere — to pay in full}

\label{lesson:LA-C419-persolvo}



\begin{quote}
Before the new one: say the Latin for interest, then the Latin for borrowed.
\end{quote}

\subsection*{You'll want to know: persolvō, persolvere}
\textbf{persolvō, persolvere} — \textquotedblleft{}to pay in full\textquotedblright{}.

\begin{quote}
Tandem omne aes aliēnum persolvit.
\end{quote}

\begin{quote}
\emph{At last he paid off his whole debt.}
\end{quote}

\emph{persolvō, persolvere, persolvī, persolūtum}, \emph{per-} through + \emph{solvō}: to loosen the whole debt.

\begin{cousinweb}
\emph{persolvō} is built from \emph{per-} \textquotedblleft{}through, completely\textquotedblright{} and \emph{solvō} \textquotedblleft{}to loosen, to pay\textquotedblright{} (Lewis \& Short).
\end{cousinweb}

\begin{grammarlens}[title={What you've built}]
One more word for telling, explaining and describing at B1.
\end{grammarlens}

\subsection*{Your turn}
\begin{itemize}
\raggedright
  \item \emph{Say it:} \emph{persolvō, persolvere}
  \item \emph{Say it:} \emph{persolvō, persolvere}, once more
  \item \emph{Say it:} a sentence of your own with \emph{persolvō, persolvere}
  \item \emph{Recall:} say the Latin for interest, then the Latin for borrowed, then say \emph{persolvō, persolvere} again
\end{itemize}

\begin{tcolorbox}[breakable,colback=teal!4,colframe=teal!35!black,title={Before you move on}]
What does persolvō, persolvere mean? (\textquotedblleft{}To pay in full\textquotedblright{}.) Say it once more.
\end{tcolorbox}
\section[ūsūra, ūsūrae]{ūsūra, ūsūrae — interest charged; use}

\label{lesson:LA-C419-usura}



\begin{quote}
Before the new one: say the Latin for borrowed, then the Latin for to pay in full.
\end{quote}

\subsection*{You'll want to know: ūsūra, ūsūrae}
\textbf{ūsūra, ūsūrae} — \textquotedblleft{}interest charged; use\textquotedblright{}.

\begin{quote}
Ūsūra gravis erat.
\end{quote}

\begin{quote}
\emph{The interest was heavy.}
\end{quote}

\emph{ūsūra, ūsūrae} f., built on \emph{ūtor}, to use: payment for the use of money, often plural. English \emph{usury}.

\begin{cousinweb}
\emph{ūsūra} is built on \emph{ūtor, ūsus} \textquotedblleft{}to use\textquotedblright{} (EtymDB; Lewis \& Short): payment for the use of money. English \emph{usury} comes from it.
\end{cousinweb}

\begin{grammarlens}[title={What you've built}]
One more word for telling, explaining and describing at B1.
\end{grammarlens}

\subsection*{Your turn}
\begin{itemize}
\raggedright
  \item \emph{Say it:} \emph{ūsūra, ūsūrae}
  \item \emph{Say it:} \emph{ūsūra, ūsūrae}, once more
  \item \emph{Say it:} a sentence of your own with \emph{ūsūra, ūsūrae}
  \item \emph{Recall:} say the Latin for borrowed, then the Latin for to pay in full, then say \emph{ūsūra, ūsūrae} again
\end{itemize}

\begin{tcolorbox}[breakable,colback=teal!4,colframe=teal!35!black,title={Before you move on}]
What does ūsūra, ūsūrae mean? (\textquotedblleft{}Interest charged; use\textquotedblright{}.) Say it once more.
\end{tcolorbox}
\section[pignus, pignoris]{pignus, pignoris — a pledge, a security}

\label{lesson:LA-C419-pignus}



\begin{quote}
Before the new one: say the Latin for to pay in full, then the Latin for interest charged.
\end{quote}

\subsection*{You'll want to know: pignus, pignoris}
\textbf{pignus, pignoris} — \textquotedblleft{}a pledge, a security\textquotedblright{}.

\begin{quote}
Ānulum prō pignore dedit.
\end{quote}

\begin{quote}
\emph{He gave a ring as a pledge.}
\end{quote}

\emph{pignus, pignoris} n. A thing handed over as security for a debt, and any pledge of love or loyalty.

\begin{cousinweb}
\emph{pignus} comes from the root of \emph{pangō} \textquotedblleft{}to fasten, to fix\textquotedblright{} (Lewis \& Short), the root also in \emph{pacīscor} \textquotedblleft{}to make a bargain\textquotedblright{}: a thing that fastens a bargain.
\end{cousinweb}

\begin{grammarlens}[title={What you've built}]
One more word for telling, explaining and describing at B1.
\end{grammarlens}

\subsection*{Your turn}
\begin{itemize}
\raggedright
  \item \emph{Say it:} \emph{pignus, pignoris}
  \item \emph{Say it:} \emph{pignus, pignoris}, once more
  \item \emph{Say it:} a sentence of your own with \emph{pignus, pignoris}
  \item \emph{Recall:} say the Latin for to pay in full, then the Latin for interest charged, then say \emph{pignus, pignoris} again
\end{itemize}

\begin{tcolorbox}[breakable,colback=teal!4,colframe=teal!35!black,title={Before you move on}]
What does pignus, pignoris mean? (\textquotedblleft{}A pledge, a security\textquotedblright{}.) Say it once more.
\end{tcolorbox}
\section[syngrapha, syngraphae]{syngrapha, syngraphae — a written bond, a promissory note}

\label{lesson:LA-C419-syngrapha}



\begin{quote}
Before the new one: say the Latin for interest charged, then the Latin for a pledge.
\end{quote}

\subsection*{You'll want to know: syngrapha, syngraphae}
\textbf{syngrapha, syngraphae} — \textquotedblleft{}a written bond, a promissory note\textquotedblright{}.

\begin{quote}
Syngrapham manū suā scrīpsit.
\end{quote}

\begin{quote}
\emph{He wrote out the bond in his own hand.}
\end{quote}

\emph{syngrapha, syngraphae} f., borrowed from Greek \emph{syngraphē}, a written contract signed by both sides.

\begin{cousinweb}
\emph{syngrapha} was borrowed from Greek \emph{syngraphē} \textquotedblleft{}a written contract\textquotedblright{} (Lewis \& Short), built from \emph{syn} \textquotedblleft{}together\textquotedblright{} and \emph{graphein} \textquotedblleft{}to write\textquotedblright{}. English \emph{syngraph} comes from it (EtymDB).
\end{cousinweb}

\begin{grammarlens}[title={What you've built}]
One more word for telling, explaining and describing at B1.
\end{grammarlens}

\subsection*{Your turn}
\begin{itemize}
\raggedright
  \item \emph{Say it:} \emph{syngrapha, syngraphae}
  \item \emph{Say it:} \emph{syngrapha, syngraphae}, once more
  \item \emph{Say it:} a sentence of your own with \emph{syngrapha, syngraphae}
  \item \emph{Recall:} say the Latin for interest charged, then the Latin for a pledge, then say \emph{syngrapha, syngraphae} again
\end{itemize}

\begin{tcolorbox}[breakable,colback=teal!4,colframe=teal!35!black,title={Before you move on}]
What does syngrapha, syngraphae mean? (\textquotedblleft{}A written bond, a promissory note\textquotedblright{}.) Say it once more.
\end{tcolorbox}
\section[fēnerātor, fēnerātōris]{fēnerātor, fēnerātōris — a moneylender}

\label{lesson:LA-C419-fenerator}



\begin{quote}
Before the new one: say the Latin for a pledge, then the Latin for a written bond.
\end{quote}

\subsection*{You'll want to know: fēnerātor, fēnerātōris}
\textbf{fēnerātor, fēnerātōris} — \textquotedblleft{}a moneylender\textquotedblright{}.

\begin{quote}
Fēnerātor ā pauperibus pecūniam repetit.
\end{quote}

\begin{quote}
\emph{The moneylender demands his money back from the poor.}
\end{quote}

\emph{fēnerātor, fēnerātōris} m., also spelled \emph{faenerātor}, built on \emph{fēnus}: a man who lends at interest.

\begin{cousinweb}
\emph{fēnerātor} is built on \emph{fēneror} \textquotedblleft{}to lend at interest\textquotedblright{} (Lewis \& Short), which comes from \emph{fēnus} \textquotedblleft{}interest\textquotedblright{}.
\end{cousinweb}

\begin{grammarlens}[title={What you've built}]
One more word for telling, explaining and describing at B1.
\end{grammarlens}

\subsection*{Your turn}
\begin{itemize}
\raggedright
  \item \emph{Say it:} \emph{fēnerātor, fēnerātōris}
  \item \emph{Say it:} \emph{fēnerātor, fēnerātōris}, once more
  \item \emph{Say it:} a sentence of your own with \emph{fēnerātor, fēnerātōris}
  \item \emph{Recall:} say the Latin for a pledge, then the Latin for a written bond, then say \emph{fēnerātor, fēnerātōris} again
\end{itemize}

\begin{tcolorbox}[breakable,colback=teal!4,colframe=teal!35!black,title={Before you move on}]
What does fēnerātor, fēnerātōris mean? (\textquotedblleft{}A moneylender\textquotedblright{}.) Say it once more.
\end{tcolorbox}
